% Options for packages loaded elsewhere
% Options for packages loaded elsewhere
\PassOptionsToPackage{unicode}{hyperref}
\PassOptionsToPackage{hyphens}{url}
\PassOptionsToPackage{dvipsnames,svgnames,x11names}{xcolor}
%
\documentclass{resolve-whitepaper}

\usepackage{xcolor}
\usepackage{amsmath,amssymb}
\setcounter{secnumdepth}{5}
\usepackage{iftex}
\ifPDFTeX
  \usepackage[T1]{fontenc}
  \usepackage[utf8]{inputenc}
  \usepackage{textcomp} % provide euro and other symbols
\else % if luatex or xetex
  \usepackage{unicode-math} % this also loads fontspec
  \defaultfontfeatures{Scale=MatchLowercase}
  \defaultfontfeatures[\rmfamily]{Ligatures=TeX,Scale=1}
\fi
\usepackage{lmodern}
\ifPDFTeX\else
  % xetex/luatex font selection
  \setmainfont[Path=C:/Users/RodrigoGordillo/AppData/Local/Temp/resolve-docs-render-4ieetd48/fonts/,Extension=.otf,UprightFont=*-Lt,BoldFont=*-Bd,ItalicFont=*-LtIt,BoldItalicFont=*-BdIt]{HelveticaNeueLTStd}
  \setsansfont[Path=C:/Users/RodrigoGordillo/AppData/Local/Temp/resolve-docs-render-4ieetd48/fonts/,Extension=.otf,UprightFont=*-Lt,BoldFont=*-Bd,ItalicFont=*-LtIt,BoldItalicFont=*-BdIt]{HelveticaNeueLTStd}
\fi
% Use upquote if available, for straight quotes in verbatim environments
\IfFileExists{upquote.sty}{\usepackage{upquote}}{}
\IfFileExists{microtype.sty}{% use microtype if available
  \usepackage[]{microtype}
  \UseMicrotypeSet[protrusion]{basicmath} % disable protrusion for tt fonts
}{}
\makeatletter
\@ifundefined{KOMAClassName}{% if non-KOMA class
  \IfFileExists{parskip.sty}{%
    \usepackage{parskip}
  }{% else
    \setlength{\parindent}{0pt}
    \setlength{\parskip}{6pt plus 2pt minus 1pt}}
}{% if KOMA class
  \KOMAoptions{parskip=half}}
\makeatother
% Make \paragraph and \subparagraph free-standing
\makeatletter
\ifx\paragraph\undefined\else
  \let\oldparagraph\paragraph
  \renewcommand{\paragraph}{
    \@ifstar
      \xxxParagraphStar
      \xxxParagraphNoStar
  }
  \newcommand{\xxxParagraphStar}[1]{\oldparagraph*{#1}\mbox{}}
  \newcommand{\xxxParagraphNoStar}[1]{\oldparagraph{#1}\mbox{}}
\fi
\ifx\subparagraph\undefined\else
  \let\oldsubparagraph\subparagraph
  \renewcommand{\subparagraph}{
    \@ifstar
      \xxxSubParagraphStar
      \xxxSubParagraphNoStar
  }
  \newcommand{\xxxSubParagraphStar}[1]{\oldsubparagraph*{#1}\mbox{}}
  \newcommand{\xxxSubParagraphNoStar}[1]{\oldsubparagraph{#1}\mbox{}}
\fi
\makeatother


\usepackage{longtable,booktabs,array}
\newcounter{none} % for unnumbered tables
\usepackage{calc} % for calculating minipage widths
% Correct order of tables after \paragraph or \subparagraph
\usepackage{etoolbox}
\makeatletter
\patchcmd\longtable{\par}{\if@noskipsec\mbox{}\fi\par}{}{}
\makeatother
% Allow footnotes in longtable head/foot
\IfFileExists{footnotehyper.sty}{\usepackage{footnotehyper}}{\usepackage{footnote}}
\makesavenoteenv{longtable}
\usepackage{graphicx}
\makeatletter
\newsavebox\pandoc@box
\newcommand*\pandocbounded[1]{% scales image to fit in text height/width
  \sbox\pandoc@box{#1}%
  \Gscale@div\@tempa{\textheight}{\dimexpr\ht\pandoc@box+\dp\pandoc@box\relax}%
  \Gscale@div\@tempb{\linewidth}{\wd\pandoc@box}%
  \ifdim\@tempb\p@<\@tempa\p@\let\@tempa\@tempb\fi% select the smaller of both
  \ifdim\@tempa\p@<\p@\scalebox{\@tempa}{\usebox\pandoc@box}%
  \else\usebox{\pandoc@box}%
  \fi%
}
% Set default figure placement to htbp
\def\fps@figure{htbp}
\makeatother





\setlength{\emergencystretch}{3em} % prevent overfull lines

\providecommand{\tightlist}{%
  \setlength{\itemsep}{0pt}\setlength{\parskip}{0pt}}



 


\makeatletter
\@ifpackageloaded{caption}{}{\usepackage{caption}}
\AtBeginDocument{%
\ifdefined\contentsname
  \renewcommand*\contentsname{Table of contents}
\else
  \newcommand\contentsname{Table of contents}
\fi
\ifdefined\listfigurename
  \renewcommand*\listfigurename{List of Figures}
\else
  \newcommand\listfigurename{List of Figures}
\fi
\ifdefined\listtablename
  \renewcommand*\listtablename{List of Tables}
\else
  \newcommand\listtablename{List of Tables}
\fi
\ifdefined\figurename
  \renewcommand*\figurename{Figure}
\else
  \newcommand\figurename{Figure}
\fi
\ifdefined\tablename
  \renewcommand*\tablename{Table}
\else
  \newcommand\tablename{Table}
\fi
}
\@ifpackageloaded{float}{}{\usepackage{float}}
\floatstyle{ruled}
\@ifundefined{c@chapter}{\newfloat{codelisting}{h}{lop}}{\newfloat{codelisting}{h}{lop}[chapter]}
\floatname{codelisting}{Listing}
\newcommand*\listoflistings{\listof{codelisting}{List of Listings}}
\makeatother
\makeatletter
\makeatother
\makeatletter
\@ifpackageloaded{caption}{}{\usepackage{caption}}
\@ifpackageloaded{subcaption}{}{\usepackage{subcaption}}
\makeatother
\usepackage{bookmark}
\IfFileExists{xurl.sty}{\usepackage{xurl}}{} % add URL line breaks if available
\urlstyle{same}
\hypersetup{
  pdftitle={ReSolve Separately Managed Account (SMA) Guide},
  colorlinks=true,
  linkcolor={brandblue},
  filecolor={Maroon},
  citecolor={brandblue},
  urlcolor={brandblue},
  pdfcreator={LaTeX via pandoc}}


\title{ReSolve Separately Managed Account (SMA) Guide}
\usepackage{etoolbox}
\makeatletter
\providecommand{\subtitle}[1]{% add subtitle to \maketitle
  \apptocmd{\@title}{\par {\large #1 \par}}{}{}
}
\makeatother
\subtitle{Understanding Your ReSolve SMA: Structure, Strategies, and
Operations}
\author{}
\date{2026-04-01}
\begin{document}
\resolveheader{ReSolve Separately Managed Account (SMA)
Guide}{Understanding Your ReSolve SMA: Structure, Strategies, and
Operations}
\makecover{ReSolve Separately Managed Account (SMA) Guide}{Understanding
Your ReSolve SMA: Structure, Strategies, and Operations}{CLIENT
GUIDE}{March 2026}
\clearpage
\tableofcontents
\clearpage

% TOC handled in partials/before-body.tex for the whitepaper format.

\section{Introduction}\label{introduction}

ReSolve Asset Management offers separately managed accounts (SMAs) for
eligible investors seeking exposure to systematic, capital-efficiency
investment strategies. An SMA allows you to access ReSolve's
institutional investment programs while maintaining full custody of your
assets. your holdings are never commingled with other clients.

This guide provides a clear overview of how ReSolve SMAs work: what
assets you can use to fund the account, how strategies are managed, what
fees apply, and what to expect on an ongoing basis. It is designed to be
a reference document throughout your relationship with ReSolve.

\sectionpagebreak

\section{How Your SMA Works}\label{how-your-sma-works}

\subsection{The Key Concept: Notional
Funding}\label{the-key-concept-notional-funding}

Unlike a traditional mutual fund or ETF, a futures-based SMA does not
require you to invest the full `portfolio size' in cash. Futures
contracts require only a fraction of their notional value as margin ---
a performance bond, not a cost of borrowing.

This means:

\begin{itemize}
\tightlist
\item
  \textbf{Your existing assets become collateral.} Cash, Treasury bills,
  stocks, bonds, ETFs, and certain mutual funds you already hold can
  serve as the margin collateral for ReSolve's futures positions. You do
  not need to liquidate your portfolio.
\item
  \textbf{Your `Nominal Account Size' is the agreed strategy
  allocation.} This is the dollar amount ReSolve manages as if it were
  the portfolio's capital base --- it is distinct from the physical
  assets in the account. The Nominal Account Size changes daily with the
  strategy's profit and loss.
\item
  \textbf{Your assets remain in your name.} All assets are held at the
  custodian in your account name. ReSolve has access to execute trades;
  it does not have the ability to withdraw assets.
\end{itemize}

\headingpagebreak

\subsection{The Two Numbers You Will
See}\label{the-two-numbers-you-will-see}

\begin{inlineexhibitblock}
\vspace{1pt}
\fontsize{8.7pt}{10.5pt}\selectfont
\setlength{\tabcolsep}{4.8pt}\renewcommand{\arraystretch}{1.0}
\renewcommand{\tabularxcolumn}[1]{m{#1}}
\arrayrulecolor{rowborder}\setlength{\arrayrulewidth}{0.25pt}
\rowcolors{2}{rowalt}{white}
\begin{tabularx}{\textwidth}{@{} >{\hsize=1.250\hsize\linewidth=\hsize\raggedright\arraybackslash\color{deepnavy}\helvMedium}X >{\hsize=0.750\hsize\linewidth=\hsize\centering\arraybackslash}X @{}}
\rowcolor{navydark}
\tableheaderstyle Concept & \tableheaderstyle What It Means \\
\tablebodystrut Nominal Account Size & The agreed portfolio size — the ‘return base.’ Changes daily with strategy P\&L. Can be increased or decreased by written instruction. \\ \hline
\tablebodystrut Investment Assets & The actual cash, T-bills, stocks, and other securities in your IBKR account serving as collateral. The amount and composition is at your discretion. \\ \hline
\end{tabularx}
\end{inlineexhibitblock}

\subsection{What You Can and Cannot
Do}\label{what-you-can-and-cannot-do}

Because the reconciliation and trade management systems at ReSolve are
automated around a single manager of the account, there are important
restrictions:

\begin{itemize}
\tightlist
\item
  \textbf{You may not place trades independently within the managed
  account.}
\item
  \textbf{You may add or withdraw collateral at any time.} Simply notify
  ReSolve's Operations team in writing using the templates provided in
  Section 5. Any additions or withdrawals do not automatically change
  your Nominal Account Size --- you must specify if you want the
  strategy allocation adjusted.
\item
  \textbf{You control what assets serve as collateral. Y}ou may hold
  cash, T-bills, equities, bonds, and eligible mutual funds as
  collateral.
\end{itemize}

\section{Account Minimums \&
Eligibility}\label{account-minimums-eligibility}

\subsection{Minimum Nominal Account
Size}\label{minimum-nominal-account-size}

\begin{inlineexhibitblock}
\vspace{1pt}
\fontsize{8.7pt}{10.5pt}\selectfont
\setlength{\tabcolsep}{4.8pt}\renewcommand{\arraystretch}{1.0}
\renewcommand{\tabularxcolumn}[1]{m{#1}}
\arrayrulecolor{rowborder}\setlength{\arrayrulewidth}{0.25pt}
\rowcolors{2}{rowalt}{white}
\begin{tabularx}{\textwidth}{@{} >{\hsize=1.250\hsize\linewidth=\hsize\raggedright\arraybackslash\color{deepnavy}\helvMedium}X >{\hsize=0.750\hsize\linewidth=\hsize\centering\arraybackslash}X @{}}
\rowcolor{navydark}
\tableheaderstyle Threshold & \tableheaderstyle Details \\
\tablebodystrut \$5,000,000 (minimum) & Accounts at this level are fully operational. Minor tracking differences relative to the model may occur due to minimum futures contract sizes, but these are symmetric and do not represent a persistent return drag. \\ \hline
\tablebodystrut \$10,000,000+ (preferred) & At this level and above, accounts closely replicate the model portfolio with negligible tracking error. \\ \hline
\end{tabularx}
\end{inlineexhibitblock}

The tracking error that can occur at the \$5M level is largely a result
of minimum futures contract sizes. Where mini-contracts are available
(the majority of markets we trade), granularity is minimal.

\subsection{Eligible Investors}\label{eligible-investors}

U.S. clients must meet applicable regulatory eligibility requirements
for managed futures accounts, including Qualified Eligible Person (QEP)
status under CFTC regulations where applicable. Please confirm your
eligibility with your advisor or ReSolve's product specialist team
before proceeding.

\subsection{Supported Brokers}\label{supported-brokers}

Interactive Brokers (IBKR), StoneX and Société Générale are supported
futures and custodial platforms available for SMA accounts.

\sectionpagebreak

\section{Eligible Collateral \& Account
Funding}\label{eligible-collateral-account-funding}

\subsection{What Can Serve as
Collateral?}\label{what-can-serve-as-collateral}

The following asset types are eligible as collateral in your IBKR
Portfolio Margin account:

\begin{inlineexhibitblock}
\vspace{1pt}
\fontsize{8.7pt}{10.5pt}\selectfont
\setlength{\tabcolsep}{4.8pt}\renewcommand{\arraystretch}{1.0}
\renewcommand{\tabularxcolumn}[1]{m{#1}}
\arrayrulecolor{rowborder}\setlength{\arrayrulewidth}{0.25pt}
\rowcolors{2}{rowalt}{white}
\begin{tabularx}{\textwidth}{@{} >{\hsize=1.250\hsize\linewidth=\hsize\raggedright\arraybackslash\color{deepnavy}\helvMedium}X >{\hsize=0.875\hsize\linewidth=\hsize\centering\arraybackslash}X >{\hsize=0.875\hsize\linewidth=\hsize\centering\arraybackslash}X @{}}
\rowcolor{navydark}
\tableheaderstyle Asset Type & \tableheaderstyle Notes & \tableheaderstyle Eligible? \\
\tablebodystrut Cash (USD) & Earns interest rate; approaches T-bill yields for larger balances & Yes \\ \hline
\tablebodystrut U.S. Treasury Bills & Preferred idle collateral; accepted for margin by IBKR & Yes \\ \hline
\tablebodystrut U.S. Treasury Bonds/Notes & Accepted with duration-based IBKR haircuts & Yes \\ \hline
\tablebodystrut Exchange-Listed Equities & Accepted subject to IBKR portfolio margin treatment & Yes* \\ \hline
\tablebodystrut ETFs & Generally accepted; IBKR margin haircuts apply & Yes \\ \hline
\tablebodystrut Mutual Funds & Eligible; note reporting nuances during settlement (see below) & Yes* \\ \hline
\tablebodystrut Illiquid / Private Assets & Not accepted as margin collateral & No \\ \hline
\end{tabularx}
\end{inlineexhibitblock}

\begin{itemize}
\tightlist
\item
  \begin{itemize}
  \tightlist
  \item
    Mutual fund purchases within the managed account create temporary
    pauses in daily reporting during the settlement period. Where
    possible, purchasing mutual funds in a separate account and
    transferring settled units into the managed account is the preferred
    approach.*
  \end{itemize}
\end{itemize}

\subsection{Adding Collateral}\label{adding-collateral}

When you wish to add assets to the account, please email Operations and
your Portfolio Manager using the following format:

\begin{callout}

Subject: Increasing collateral value in my account

\end{callout}

I will be adding \${[}AMOUNT{]} of collateral in the form of {[}cash /
T-bills / stock ticker / mutual fund ticker{]} on {[}DATE{]}. Please
maintain existing nominal account size. Send to:
operations@investresolve.com --- CC your Portfolio Manager

\begin{callout}

\end{callout}

\subsection{Withdrawing Collateral}\label{withdrawing-collateral}

\begin{callout}

Subject: Decreasing collateral value in my account

\end{callout}

I will be withdrawing \${[}AMOUNT{]} of collateral in the form of
{[}cash / T-bills / stock ticker / mutual fund ticker{]} on {[}DATE{]}.
Please maintain existing nominal account size. Send to:
operations@investresolve.com --- CC your Portfolio Manager

\begin{callout}

\end{callout}

\section{Investment Strategies}\label{investment-strategies}

ReSolve offers the following strategies in SMA format. A key design
feature of each strategy is that the target volatility level is a
configurable parameter --- agreed upon with you at account setup ---
rather than a fixed attribute. This allows us to tailor the strategy's
risk profile to your objectives.

\subsection{Systematic Macro}\label{systematic-macro}

\textbf{What it does:} Seeks to generate consistent capital appreciation
by deploying a diversified ensemble of systematic alpha strategies
across global exchange-traded futures markets. The strategy harvests a
variety of return drivers including trend, momentum, carry, relative
value, volatility, skewness, and seasonality.

\textbf{How it works:} Advanced statistical and machine learning methods
identify marginal market inefficiencies across a broad universe of
global futures. Positions are sized to maximize risk-adjusted return.
Portfolio exposure is actively managed to maintain the target volatility
level --- scaling up in calm markets and scaling down during periods of
elevated volatility and correlation.

\textbf{Volatility target:} Configurable. Common range: 12\%--20\%
annualized. A 16\% target is a frequently used default.

\textbf{Instruments:} Global exchange-traded futures across equities,
fixed income, currencies, commodities, and other markets.

\textbf{Fee structure:} Management fee + Performance fee (see Section
6).

\subsection{Futures Yield Program
(Carry)}\label{futures-yield-program-carry}

\textbf{What it does:} Systematically harvests the carry premium
embedded in global futures markets --- capturing returns from holding
positions where forward prices are at a discount to spot (backwardation)
or managing the cost in contango markets.

\textbf{How it works:} The expanded-universe version of this strategy
trades a broad set of global futures markets, identifying and sizing
carry opportunities systematically with ongoing risk management to
control drawdowns.

\textbf{Volatility target:} Configurable. A 20\% target is commonly
used. A smaller-universe version of this strategy is also available in
ETF format for investors who prefer a liquid fund structure.

\textbf{Fee structure:} Management fee only (no performance fee).

\subsection{All Terrain}\label{all-terrain}

\textbf{What it does:} Delivers diversified risk-adjusted returns across
all market environments by combining three complementary sleeves: a
risk-balanced beta portfolio, a fast-moving systematic macro overlay,
and a long-volatility sleeve for tail-risk protection.

\textbf{How it works:} The beta sleeve holds a risk-balanced allocation
across global equities (run as a whole, not split by region), bonds,
gold, and commodities --- managed using long-term trend and
mean-reversion signals. The systematic macro overlay deploys quickly
during market corrections to dampen beta losses. The long-volatility
sleeve is designed to trigger specifically during abrupt equity market
shocks (when VIX term structure inverts and implied volatility exceeds
realized volatility), typically associated with equity drawdowns of
10--15\%+.

\textbf{Volatility target:} {[}TBD --- confirm for SMA format.{]}

\textbf{Availability:} Available in both SMA format and U.S. fund
structure.

\textbf{Fee structure:} Management fee + Performance fee with a hurdle
rate (see Section 6).

\subsection{Trend Replication}\label{trend-replication}

\textbf{What it does:} Captures trend-following returns across global
futures markets using a disciplined, rules-based replication approach.

\textbf{Volatility target:} {[}TBD{]}.

\textbf{Fee structure:} Management fee only (no performance fee).

\sectionpagebreak

\section{Volatility Targeting \& Risk}\label{volatility-targeting-risk}

\subsection{What `Volatility Target'
Means}\label{what-volatility-target-means}

The volatility target is the annualized standard deviation of returns
the strategy aims to maintain. It is a risk measure --- not a return
guarantee. The strategy's actual volatility may be higher or lower than
the target depending on market conditions.

To maintain the target, ReSolve actively adjusts the strategy's gross
exposure to futures markets:

\begin{itemize}
\tightlist
\item
  When market volatility is low and correlations are low, the strategy
  may hold more exposure to hit its target.
\item
  When market volatility is high and correlations spike, the strategy
  reduces exposure to prevent realized volatility from exceeding the
  target.
\end{itemize}

\subsection{Leverage and Risk}\label{leverage-and-risk}

\textbf{Important:} Higher volatility targets imply higher gross
exposure (leverage). Leveraged portfolios can lose more than the initial
collateral deposited. The Investment Guidelines you sign will include an
explicit acknowledgment of this risk.

\begin{inlineexhibitblock}
\vspace{1pt}
\fontsize{8.7pt}{10.5pt}\selectfont
\setlength{\tabcolsep}{4.8pt}\renewcommand{\arraystretch}{1.0}
\renewcommand{\tabularxcolumn}[1]{m{#1}}
\arrayrulecolor{rowborder}\setlength{\arrayrulewidth}{0.25pt}
\rowcolors{2}{rowalt}{white}
\begin{tabularx}{\textwidth}{@{} >{\hsize=1.250\hsize\linewidth=\hsize\raggedright\arraybackslash\color{deepnavy}\helvMedium}X >{\hsize=0.875\hsize\linewidth=\hsize\centering\arraybackslash}X >{\hsize=0.875\hsize\linewidth=\hsize\centering\arraybackslash}X @{}}
\rowcolor{navydark}
\tableheaderstyle Target Volatility & \tableheaderstyle Risk Profile & \tableheaderstyle Typical Use Case \\
\tablebodystrut 10–12\% & Moderate & Core diversifier, lower drawdown tolerance \\ \hline
\tablebodystrut 14–16\% & Moderate-High & Balanced growth and diversification \\ \hline
\tablebodystrut 18–20\%+ & High & Satellite allocation, maximum return potential \\ \hline
\end{tabularx}
\end{inlineexhibitblock}

\sectionpagebreak

\section{Fee Structure}\label{fee-structure}

\subsection{Overview by Strategy}\label{overview-by-strategy}

\begin{inlineexhibitblock}
\vspace{1pt}
\fontsize{8.7pt}{10.5pt}\selectfont
\setlength{\tabcolsep}{4.8pt}\renewcommand{\arraystretch}{1.0}
\renewcommand{\tabularxcolumn}[1]{m{#1}}
\arrayrulecolor{rowborder}\setlength{\arrayrulewidth}{0.25pt}
\rowcolors{2}{rowalt}{white}
\begin{tabularx}{\textwidth}{@{} >{\hsize=1.250\hsize\linewidth=\hsize\raggedright\arraybackslash\color{deepnavy}\helvMedium}X >{\hsize=0.917\hsize\linewidth=\hsize\centering\arraybackslash}X >{\hsize=0.917\hsize\linewidth=\hsize\centering\arraybackslash}X >{\hsize=0.917\hsize\linewidth=\hsize\centering\arraybackslash}X @{}}
\rowcolor{navydark}
\tableheaderstyle Alpha Strategies & \tableheaderstyle Management Fee & \tableheaderstyle Performance Fee & \tableheaderstyle Performance Fee Billing \\
\tablebodystrut Systematic Macro & 1.50\% per annum & 15\% above High-Water Mark & Quarterly \\ \hline
\tablebodystrut All Terrain & 1.00\% per annum* & 10\% above Benchmark* above a High-Water Mark & Quarterly \\ \hline
\end{tabularx}
\tabnote{All Terrain fee reflects the Founders Class. Confirm current availability and hurdle rate benchmark with your Portfolio Manager.*}
\end{inlineexhibitblock}

\begin{inlineexhibitblock}
\vspace{1pt}
\fontsize{8.7pt}{10.5pt}\selectfont
\setlength{\tabcolsep}{4.8pt}\renewcommand{\arraystretch}{1.0}
\renewcommand{\tabularxcolumn}[1]{m{#1}}
\arrayrulecolor{rowborder}\setlength{\arrayrulewidth}{0.25pt}
\rowcolors{2}{rowalt}{white}
\begin{tabularx}{\textwidth}{@{} >{\hsize=1.250\hsize\linewidth=\hsize\raggedright\arraybackslash\color{deepnavy}\helvMedium}X >{\hsize=0.875\hsize\linewidth=\hsize\centering\arraybackslash}X >{\hsize=0.875\hsize\linewidth=\hsize\centering\arraybackslash}X @{}}
\rowcolor{navydark}
\tableheaderstyle Factor Strategies & \tableheaderstyle Management Fee & \tableheaderstyle Performance Fee \\
\tablebodystrut Futures Yield / Carry & 0.95\% per annum & None \\ \hline
\tablebodystrut Trend Replication & 0.95\% per annum & None \\ \hline
\end{tabularx}
\end{inlineexhibitblock}

\subsection{Management Fee}\label{management-fee}

The management fee accrues daily, calculated on that day's Nominal
Account Size. At the end of each month, the total daily accruals are
deducted from the account.

\subsection{Performance Fee \& High-Water
Mark}\label{performance-fee-high-water-mark}

\begin{itemize}
\tightlist
\item
  \textbf{What triggers the fee:} A performance fee is charged only when
  the quarter-end Nominal Account Size exceeds the previous High-Water
  Mark --- meaning the strategy has generated new net profits above the
  previous peak.
\item
  \textbf{High-Water Mark (HWM):} Set at the Nominal Account Size at the
  end of the most recent quarter in which a performance fee was paid. If
  no fee has been paid since account opening, the HWM is the initial
  Nominal Account Size, adjusted proportionately for any withdrawals.
\item
  \textbf{No double-charging:} Performance fees are only charged on new
  gains above the HWM. Prior losses must be recovered before any further
  performance fee accrues.
\item
  \textbf{Billing cycle:} Performance fees are calculated and charged on
  a calendar quarter basis.
\end{itemize}

\section{Cash Management}\label{cash-management}

\subsection{How ReSolve Manages Cash in Your Account (IBKR
Example)}\label{how-resolve-manages-cash-in-your-account-ibkr-example}

Because futures trading generates daily P\&L and variation margin flows,
your account will typically hold some idle cash at any point during the
month. ReSolve manages this cash to optimize returns while maintaining
the liquidity needed for ongoing margin requirements.

\subsection{Standard Protocol}\label{standard-protocol}

\begin{itemize}
\tightlist
\item
  \textbf{Monthly sweep:} At month-end, excess cash above an operational
  buffer is swept into short-dated U.S. Treasury bills.
\item
  \textbf{10\% intra-month threshold:} If cash balances grow to exceed
  10\% of the Nominal Account Size intra-month, ReSolve will purchase
  T-bills before the scheduled month-end sweep.
\item
  \textbf{Buffer maintenance:} A cash buffer is always maintained to
  cover margin fluctuations and operational needs. The size of this
  buffer is set by the trading team based on current market conditions.
\end{itemize}

\subsection{Cash vs.~T-Bills: The
Economics}\label{cash-vs.-t-bills-the-economics}

For most clients, the difference between holding cash and T-bills is
small but worth understanding:

\begin{itemize}
\tightlist
\item
  IBKR pays interest on USD cash balances, but only above a minimum
  threshold and typically at a rate slightly below prevailing T-bill
  yields.
\item
  T-bills earn their own yield directly.
\item
  If the account holds a negative cash balance (i.e., borrowing against
  T-bill collateral for foreign margin), IBKR charges a borrowing rate
  typically 0.5--1.0\% above T-bill yields. The after-tax cost of this
  is often modest, given the deductibility of margin interest.
\end{itemize}

\section{Margin \& the Nominal Account
Size}\label{margin-the-nominal-account-size}

\subsection{How Margin Works in Your
Account}\label{how-margin-works-in-your-account}

Futures margin is not a borrowing cost --- it is a performance bond
posted with the exchange to guarantee contract obligations. The actual
margin consumed by ReSolve's futures positions is a fraction of the
Nominal Account Size. Your Investment Assets (T-bills, stocks, bonds,
etc.) collectively satisfy this margin requirement.

IBKR's Portfolio Margin framework applies margin requirements based on
the net risk of the overall portfolio, enabling efficient capital
utilization across diversified, uncorrelated futures positions.

\subsection{What Happens if Margin Is
Insufficient?}\label{what-happens-if-margin-is-insufficient}

If the value of your Investment Assets falls below the level needed to
support the existing futures positions, ReSolve will take one of two
steps:

\begin{itemize}
\tightlist
\item
  \textbf{Request additional collateral:} You will be notified in
  writing and asked to deposit additional assets.
\item
  \textbf{Reduce the Nominal Account Size:} If additional collateral is
  not provided, ReSolve may reduce the Nominal Account Size to a level
  consistent with available margin. The Nominal Account Size will be
  restored as permitted by improved conditions. Any such change will be
  communicated to you in writing.
\end{itemize}

\subsection{Changing Your Nominal Account
Size}\label{changing-your-nominal-account-size}

You can request an increase or decrease to your Nominal Account Size at
any time by written instruction to your Portfolio Manager and
Operations. Note that simply adding or removing collateral without an
accompanying instruction does not change the Nominal Account Size ---
the strategy continues running at its current authorized level.

\sectionpagebreak

\section{Reporting \& Communications}\label{reporting-communications}

\subsection{Daily Performance Report}\label{daily-performance-report}

You will have full access to your custodial account including detailed
account statements, showing all positions, cash flows, interest credits
and debits, and transaction history.

ReSolve will provide a quarterly performance report at the end of each
calendar quarter.

\subsection{How to Contact Us}\label{how-to-contact-us}

\begin{inlineexhibitblock}
\vspace{1pt}
\fontsize{8.7pt}{10.5pt}\selectfont
\setlength{\tabcolsep}{4.8pt}\renewcommand{\arraystretch}{1.0}
\renewcommand{\tabularxcolumn}[1]{m{#1}}
\arrayrulecolor{rowborder}\setlength{\arrayrulewidth}{0.25pt}
\rowcolors{2}{rowalt}{white}
\begin{tabularx}{\textwidth}{@{} >{\hsize=1.250\hsize\linewidth=\hsize\raggedright\arraybackslash\color{deepnavy}\helvMedium}X >{\hsize=0.875\hsize\linewidth=\hsize\centering\arraybackslash}X >{\hsize=0.875\hsize\linewidth=\hsize\centering\arraybackslash}X @{}}
\rowcolor{navydark}
\tableheaderstyle Team & \tableheaderstyle Email & \tableheaderstyle For \\
\tablebodystrut Operations & operations@investresolve.com & Collateral changes, account setup, trade confirmations \\ \hline
\tablebodystrut Portfolio Management & Via your Portfolio Manager & Strategy questions, Nominal Account Size changes, performance \\ \hline
\end{tabularx}
\end{inlineexhibitblock}

\section{Tax Considerations}\label{tax-considerations}

\textbf{Disclaimer:} \emph{The following overview is for informational
purposes only. It does not constitute tax advice. Please consult your
tax advisor regarding your specific circumstances.}

\subsection{Section 1256 Contracts: The 60/40
Rule}\label{section-1256-contracts-the-6040-rule}

Most U.S.-listed exchange-traded futures contracts qualify as Section
1256 contracts under the Internal Revenue Code. These receive a blended
capital gains tax treatment:

\begin{itemize}
\tightlist
\item
  60\% of net gains are treated as long-term capital gains (lower tax
  rate).
\item
  40\% of net gains are treated as short-term capital gains (ordinary
  income rates).
\end{itemize}

This 60/40 treatment applies regardless of how long positions are held.
It is one of the most tax-efficient features of futures-based investment
strategies compared to short-term equity trading.

Additionally, Section 1256 losses can be carried back up to three years,
providing a unique tax planning tool not available with most other
capital losses.

\subsection{Non-Section 1256 Contracts (Foreign
Futures)}\label{non-section-1256-contracts-foreign-futures}

Some foreign futures contracts do not qualify as Section 1256 contracts.
For these positions:

\begin{itemize}
\tightlist
\item
  Tax treatment is based on actual holding period. Most foreign futures
  positions are held for shorter durations and are likely to generate
  short-term capital gains taxed at ordinary income rates.
\item
  The average holding period for non-Section 1256 positions in the
  Systematic Macro strategy has historically been approximately 90 days,
  though this varies with market conditions.
\end{itemize}

\subsection{Interest Income and
Expense}\label{interest-income-and-expense}

\begin{itemize}
\tightlist
\item
  Interest earned on T-bills and cash is generally taxable as ordinary
  income at the federal level and may be exempt from state income tax
  (for U.S. Treasury obligations).
\item
  Interest paid on margin borrowing is generally deductible as
  investment interest expense, subject to applicable limitations.
\end{itemize}

\sectionpagebreak

\section{Getting Started: Onboarding
Overview}\label{getting-started-onboarding-overview}

\subsection{The Process}\label{the-process}

\begin{inlineexhibitblock}
\vspace{1pt}
\fontsize{8.7pt}{10.5pt}\selectfont
\setlength{\tabcolsep}{4.8pt}\renewcommand{\arraystretch}{1.0}
\renewcommand{\tabularxcolumn}[1]{m{#1}}
\arrayrulecolor{rowborder}\setlength{\arrayrulewidth}{0.25pt}
\rowcolors{2}{rowalt}{white}
\begin{tabularx}{\textwidth}{@{} >{\hsize=0.400\hsize\linewidth=\hsize\raggedright\arraybackslash\color{deepnavy}\helvMedium}X >{\hsize=0.800\hsize\linewidth=\hsize\centering\arraybackslash}X >{\hsize=1.800\hsize\linewidth=\hsize\centering\arraybackslash}X @{}}
\rowcolor{navydark}
\tableheaderstyle Step & \tableheaderstyle Stage & \tableheaderstyle What Happens \\
\tablebodystrut 1 & Initial Discussion & Agree on strategy, volatility target, initial Nominal Account Size, and fee structure with your Portfolio Manager. \\ \hline
\tablebodystrut 2 & Documentation & Sign the Investment Guidelines Agreement (specifying strategy, Nominal Account Size, volatility target, fee schedule, and IBKR account number). \\ \hline
\tablebodystrut 3 & Account Setup & Open or confirm IBKR Portfolio Margin account; grant ReSolve Limited Power of Attorney for trade execution. \\ \hline
\tablebodystrut 4 & Funding & Transfer your chosen collateral assets into the IBKR account. Notify Operations per the collateral addition template. \\ \hline
\tablebodystrut 5 & Strategy Launch & ReSolve establishes futures positions at the agreed Nominal Account Size. Daily reporting begins once account is in steady state. \\ \hline
\tablebodystrut 6 & Ongoing & Receive daily performance reports. ReSolve manages cash sweeps, margin, and rebalancing. Review monthly IBKR statements. \\ \hline
\end{tabularx}
\end{inlineexhibitblock}

\section{Frequently Asked Questions}\label{frequently-asked-questions}

\subsection{Can I hold my existing stock and bond portfolio as
collateral?}\label{can-i-hold-my-existing-stock-and-bond-portfolio-as-collateral}

Yes. You do not need to liquidate your current portfolio to fund the
SMA. Eligible assets (equities, bonds, ETFs, T-bills, and certain mutual
funds) can serve as collateral within the IBKR portfolio margin
framework. IBKR applies haircuts to non-cash collateral, so heavily
concentrated or volatile equity positions may contribute less to margin
capacity than their market value suggests.

\subsection{What happens to my account if markets sell off
sharply?}\label{what-happens-to-my-account-if-markets-sell-off-sharply}

ReSolve's volatility-targeting approach naturally reduces futures
exposure during periods of elevated market volatility, providing some
built-in downside mitigation. However, no approach eliminates risk
entirely, and leveraged futures strategies can sustain significant
losses. If losses reduce Investment Assets to the point where margin is
insufficient, ReSolve will request additional collateral or reduce the
Nominal Account Size to a sustainable level.

\subsection{Why does my account sometimes show a large cash
balance?}\label{why-does-my-account-sometimes-show-a-large-cash-balance}

Cash accumulates intra-month from daily trading P\&L, margin
fluctuations, and timing of collateral movements. ReSolve sweeps excess
cash into T-bills at month-end (and sooner if cash exceeds 10\% of the
Nominal Account Size). Temporary cash balances do not represent an
inefficiency --- they are a normal part of how futures-based accounts
operate.

\subsection{Can I add to the account
mid-month?}\label{can-i-add-to-the-account-mid-month}

Yes. Simply notify Operations and your Portfolio Manager using the
standard collateral addition template. Specify whether you want the
Nominal Account Size increased (i.e., the strategy allocation expanded)
or simply the collateral buffer increased without changing the strategy
scale.

\subsection{Will the strategy be affected by my other IBKR
accounts?}\label{will-the-strategy-be-affected-by-my-other-ibkr-accounts}

No.~The SMA is managed in a single IBKR account that ReSolve operates
independently. Your other IBKR accounts (or accounts at other brokers)
are separate and have no impact on the managed account.

\subsection{Can my account be managed in Canada in the
future?}\label{can-my-account-be-managed-in-canada-in-the-future}

Canada and other jurisdictions are on ReSolve's roadmap for SMA
availability. Please speak with your Portfolio Manager for the latest
timeline.

\subsection{What are the tax benefits of futures-based
strategies?}\label{what-are-the-tax-benefits-of-futures-based-strategies}

U.S. exchange-traded futures contracts benefit from Section 1256 tax
treatment --- 60\% long-term / 40\% short-term capital gains regardless
of holding period. Section 1256 losses can also be carried back up to
three years. These features make futures-based strategies comparatively
tax-efficient. Consult your tax advisor for guidance on your specific
situation.

\headingpagebreak

\disclosuresection{Important Disclosures}

\begin{disclosurebody}

ReSolve Asset Management is registered in the United States with the
Commodity Futures Trading Commission (CFTC) as a Commodity Trading
Advisor (CTA) and Commodity Pool Operator (CPO), administered through
the National Futures Association (NFA). ReSolve is also registered with
the Securities and Exchange Commission (SEC) as a Registered Investment
Advisor (RIA).

Past performance is not indicative of future results. The strategies
described in this guide use leverage and can exhibit significant
volatility. The use of leverage can lead to substantial portfolio
losses, including losses that exceed the initial investment. Separately
managed accounts are not appropriate for all investors.

This guide is intended as a general reference and does not constitute
investment, legal, or tax advice. Clients should review their Investment
Guidelines Agreement for the specific terms governing their account. Any
discrepancies between this guide and the Investment Guidelines Agreement
are governed by the Agreement.

\emph{Version 1.0 (Draft) \textbar{} March 2026 \textbar{} ReSolve Asset
Management}

\end{disclosurebody}




\end{document}
